\documentclass[11pt,a4paper]{article}
\usepackage{fontspec}
\IfFontExistsTF{DejaVu Serif}{\setmainfont{DejaVu Serif}}{\setmainfont{Times New Roman}}
\IfFontExistsTF{DejaVu Sans Mono}{\setmonofont[Scale=0.85]{DejaVu Sans Mono}}{\setmonofont[Scale=0.85]{Courier New}}
\usepackage[russian]{babel}
\usepackage[margin=19mm]{geometry}
\usepackage{amsmath,amssymb,booktabs,array,longtable}
\usepackage{xcolor}
\usepackage[unicode,colorlinks=true,linkcolor=blue,urlcolor=blue]{hyperref}
\setlength{\parindent}{0pt}
\setlength{\parskip}{5pt}
\setlength{\emergencystretch}{5em}
\renewcommand{\arraystretch}{1.16}
\newcommand{\X}{BS^2}
\newcommand{\bytes}{\text{ bytes}}
\newcommand{\page}[1]{\clearpage\section{#1}}
\newcommand{\note}[1]{\par\textbf{Пояснение.} #1\par}
\begin{document}
\begin{center}
{\LARGE HW1, версия 2: полный вывод формул}\\[5pt]
{\large Основа для самостоятельной рукописной работы}\\[3pt]
\small По \texttt{models.py}, \texttt{equations.py}, \texttt{calibrate.py}\\
и прогону \texttt{20260927\_090211\_ab1cd7d2}. 27 сентября 2026 г.
\end{center}

Здесь разобрана именно последовательная \texttt{SmallCNN} из HW1: шесть
свёрток, без residual-связей. Версия 2 относится к расчётам и калибровке;
архитектура сети та же. Текст можно переписывать по разделам: объяснение,
формула, подстановка, результат. Исходный рукописный PDF сохраняется отдельно.

\section{Обозначения и размеры тензоров}
$S$ --- сторона квадратного изображения, $B$ --- число изображений в batch.
В задании $S\ge32$ кратно 16, $B$ --- положительное целое. Для сокращения
записи положим $X=BS^2$. Вход имеет форму $B\times3\times S\times S$ и
содержит $3X$ чисел. FP32 занимает 4 bytes; $1\,\mathrm{MiB}=2^{20}$ bytes.

Работаем в \texttt{eval()} и \texttt{inference\_mode()}, без градиентов.
Все Conv имеют \texttt{bias=False}, $p=\lfloor k/2\rfloor$, dilation 1;
после каждой --- BatchNorm и ReLU \texttt{inplace=True}.
\texttt{cudnn.benchmark=False}; оба флага \texttt{allow\_tf32=False}.

Для входной стороны $h$, ядра $k$, padding $p$ и шага $s$:
\[
 h_{out}=\left\lfloor\frac{h+2p-k}{s}\right\rfloor+1.
\]
Это число допустимых позиций ядра: стартуем с 0, смещаемся на $s$,
последняя позиция не превосходит $h+2p-k$. Для Conv1:
\[
\left\lfloor\frac{S+6-7}{2}\right\rfloor+1
=\left\lfloor\frac{S-1}{2}\right\rfloor+1=S/2.
\]
Последнее равенство верно для чётного $S$; округление вниз нельзя терять.
MaxPool далее даёт $S/4$, следующие stride-2 Conv --- $S/8$ и $S/16$.

\begin{center}\small
\begin{tabular}{lrrrrl}
\toprule
Операция & $C_{in}$ & $C_{out}$ & $k$ & $s$ & Выход / число FP32 \\
\midrule
Conv1--BN1--ReLU & 3 & 32 & 7 & 2 & $B\times32\times(S/2)^2$: $8X$ \\
MaxPool, $p=1$ & 32 & 32 & 3 & 2 & $B\times32\times(S/4)^2$: $2X$ \\
Conv2--BN2--ReLU & 32 & 64 & 5 & 1 & $B\times64\times(S/4)^2$: $4X$ \\
Conv3--BN3--ReLU & 64 & 128 & 3 & 2 & $B\times128\times(S/8)^2$: $2X$ \\
Conv4--BN4--ReLU & 128 & 256 & 1 & 1 & $B\times256\times(S/8)^2$: $4X$ \\
Conv5--BN5--ReLU & 256 & 256 & 3 & 2 & $B\times256\times(S/16)^2$: $X$ \\
Conv6--BN6--ReLU & 256 & 512 & 1 & 1 & $B\times512\times(S/16)^2$: $2X$ \\
GlobalAvgPool & 512 & 512 & --- & --- & $B\times512\times1\times1$: $512B$ \\
Flatten & --- & --- & --- & --- & $B\times512$: $512B$, view \\
Linear1--ReLU & 512 & 256 & --- & --- & $B\times256$: $256B$ \\
Linear2 & 256 & 100 & --- & --- & $B\times100$: $100B$ \\
\bottomrule
\end{tabular}\end{center}
BN и ReLU не меняют форму. В Linear есть bias. Flatten в данной сети меняет
представление формы без нового массива данных.


\page{GlobalAvgPool и Flatten: как получается размер выхода}
После последнего блока Conv6--BN6--ReLU имеем тензор
\[
 A\in\mathbb R^{B\times512\times h\times h},\qquad h=S/16.
\]
Здесь $A_{b,c,u,v}$ --- число в изображении $b$, канале $c$, строке $u$
и столбце $v$. Для каждого изображения существует 512 отдельных карт
признаков размером $h\times h$.

\textbf{GlobalAvgPool усредняет пространственные позиции одного канала.}
В коде это \texttt{AdaptiveAvgPool2d(1)}: требуемый размер выхода по высоте
и ширине явно задан как $1\times1$. Для фиксированных $b,c$:
\[
 g_{b,c,1,1}=\frac{1}{h^2}\sum_{u=1}^{h}\sum_{v=1}^{h}A_{b,c,u,v}.
\]
Сначала складываем все $h^2$ чисел карты, затем делим на их количество.
Например, карта $2\times2$ превращается в
$(A_{b,c,1,1}+A_{b,c,1,2}+A_{b,c,2,1}+A_{b,c,2,2})/4$.
Каналы между собой не смешиваются, разные изображения также не смешиваются.
Поэтому $B$ и 512 сохраняются, а две пространственные оси становятся единицами:
\[
 B\times512\times(S/16)\times(S/16)
 \ \longrightarrow\ B\times512\times1\times1.
\]
Число значений уменьшается с $B\cdot512\cdot(S/16)^2=2BS^2$
до $B\cdot512\cdot1\cdot1=512B$.
Это не одна средняя величина на изображение: средних величин \emph{512}.
У GAP нет обучаемых весов или bias.

\textbf{Арифметика.} Для одной карты из $n=h^2$ чисел нужны $n-1$ сложений
и одно деление. Карт $512B$, поэтому по принятому счёту операций:
\[
 F_{GAP}=512B\bigl[(h^2-1)+1\bigr]
 =512Bh^2=512B\frac{S^2}{256}=2BS^2.
\]
В нашей модели трафика читаем весь вход и записываем все средние:
\[
 Q_{GAP}=4(2BS^2+512B)\ \text{bytes}.
\]
Сам результат GAP занимает $4\cdot512B=2048B$ bytes FP32-данных
до округления allocator; это размер результата, не пик всей сети.

\textbf{Flatten объединяет оси признаков, оставляя batch отдельно.}
В \texttt{Flatten()} по умолчанию объединяются оси начиная с первой после
batch. Произведение их размеров $512\cdot1\cdot1=512$:
\[
 B\times512\times1\times1\ \longrightarrow\ B\times512,
 \qquad z_{b,c}=g_{b,c,1,1}.
\]
Значения не меняются, усреднения и умножения нет. В данной сети это view
того же storage: $F_{Flatten}=0$, модельный трафик данных $Q_{Flatten}=0$,
новый массив из $512B$ чисел не выделяется.

\textbf{Пример.} При $S=224$, $B=8$ имеем $h=14$, $h^2=196$.
GAP превращает $8\times512\times14\times14$ в $8\times512\times1\times1$:
для каждого из $8\cdot512=4096$ результатов выполняет 195 сложений и деление.
Итого $4096\cdot196=802816$ FLOPs. Flatten даёт матрицу $8\times512$
из тех же 4096 значений.

\page{Linear и ReLU: от 512 признаков к 100 выходам}
После Flatten каждое изображение описывается вектором из 512 признаков.
Матрица входов $Z=(z_{b,c})$ имеет размер $B\times512$.
Размеры 256 и 100 заданы архитектурой: это числа выходов соответствующих
Linear, а не результат формулы stride/padding.

\textbf{Linear1: \texttt{Linear(512,256)}.} Есть матрица весов
$W^{(1)}\in\mathbb R^{256\times512}$ и вектор bias
$\beta^{(1)}\in\mathbb R^{256}$. Для изображения $b$ и выхода $j$:
\[
 u_{b,j}=\sum_{c=1}^{512}W^{(1)}_{j,c}\,z_{b,c}+\beta^{(1)}_j,
 \qquad j=1,\ldots,256.
\]
Каждый выход использует все 512 входных признаков со своими весами.
В матричной форме $U=Z(W^{(1)})^T+\beta^{(1)}$, где bias прибавляется
к каждой строке. Формы матриц:
\[
 (B\times512)(512\times256)\ \longrightarrow\ B\times256.
\]
Выход содержит $256B$ чисел. Параметров $256\cdot512+256=131328$,
они общие для всех изображений и не умножаются на $B$.
По соглашению $1\,MAC=2\,FLOPs$ и с отдельным сложением bias:
\[
 F_{L1}=B\cdot256\,(2\cdot512+1)=262400B.
\]

\textbf{ReLU после Linear1.} Для каждого элемента:
\[
 r_{b,j}=\max(0,u_{b,j}),\qquad R\in\mathbb R^{B\times256}.
\]
Отрицательные значения заменяются нулём, остальные сохраняются.
Форма $B\times256$ не меняется. Сравнения исключены из нашего FLOPs,
поэтому $F_{ReLU}=0$. При inplace новый массив не создаётся, но чтение
и запись остаются: $Q_{ReLU}=4\cdot2\cdot256B=2048B$ bytes в модели трафика.

\textbf{Linear2: \texttt{Linear(256,100)}.} Веса
$W^{(2)}\in\mathbb R^{100\times256}$, bias $\beta^{(2)}\in\mathbb R^{100}$:
\[
 y_{b,\ell}=\sum_{j=1}^{256}W^{(2)}_{\ell,j}\,r_{b,j}+\beta^{(2)}_\ell,
 \qquad \ell=1,\ldots,100.
\]
Матрица выхода $Y=R(W^{(2)})^T+\beta^{(2)}$ имеет размер $B\times100$.
Всего $100B$ выходных чисел и $100\cdot256+100=25700$ параметров.
\[
 F_{L2}=B\cdot100\,(2\cdot256+1)=51300B.
\]
Это 100 оценок классов (logits) на изображение; Softmax в данной сети нет.
Во всех формулах bias $\beta$ отличается от обозначения batch size $B$.

\textbf{Численная цепочка при $B=8$.}
$8\times512\to8\times256\to8\times256\to8\times100$.
Linear1: $262400\cdot8=2099200$ FLOPs; Linear2:
$51300\cdot8=410400$ FLOPs. При $S=224$ весь head вместе с GAP даёт
$802816+2099200+410400=3312416$ FLOPs по нашему соглашению MAC.

\page{FLOPs: от одного выходного числа ко всем свёрткам}
Для одного выходного числа Conv используется $C_{in}k^2$ произведений
входов на веса. По соглашению задания один multiply--accumulate (MAC)
считается как 2 FLOPs. Поэтому стоимость одного выходного числа равна
$2C_{in}k^2$, а таких чисел $BC_{out}h_{out}^2$:
\[
\boxed{F_{Conv}=2BC_{out}h_{out}^2C_{in}k^2.}
\]
Здесь $B$ учитывает все изображения, $C_{out}$ --- все фильтры,
$h_{out}^2$ --- все позиции. Веса одинаковы для разных позиций и изображений,
но арифметические операции повторяются. Padding не вычитаем из стандартного
плотного подсчёта: считаем и позиции с дополненными нулями.

\textbf{Conv1: $7\times7$, $3\to32$, выход $S/2$.}
\[
 C_{in}k^2=3\cdot49=147,\quad N_{out}=8X,
 \quad F_1=2\cdot147\cdot8X=2352X.
\]
\textbf{Conv2: $5\times5$, $32\to64$, выход $S/4$.}
\[
 C_{in}k^2=32\cdot25=800,\quad N_{out}=4X,
 \quad F_2=2\cdot800\cdot4X=6400X.
\]
\textbf{Conv3: $3\times3$, $64\to128$, выход $S/8$.}
\[
 C_{in}k^2=64\cdot9=576,\quad N_{out}=2X,
 \quad F_3=2\cdot576\cdot2X=2304X.
\]
\textbf{Conv4: $1\times1$, $128\to256$, выход $S/8$.}
\[
 C_{in}k^2=128,\quad N_{out}=4X,
 \quad F_4=2\cdot128\cdot4X=1024X.
\]
\textbf{Conv5: $3\times3$, $256\to256$, выход $S/16$.}
\[
 C_{in}k^2=256\cdot9=2304,\quad N_{out}=X,
 \quad F_5=2\cdot2304\cdot X=4608X.
\]
\textbf{Conv6: $1\times1$, $256\to512$, выход $S/16$.}
\[
 C_{in}k^2=256,\quad N_{out}=2X,
 \quad F_6=2\cdot256\cdot2X=1024X.
\]
Суммируем \emph{все} свёртки:
\[
\boxed{F_{Conv,total}=(2352+6400+2304+1024+4608+1024)X=17712X.}
\]

\note{Это точный результат принятого соглашения $1\,MAC=2\,FLOPs$,
а не число машинных инструкций GPU. Если считать скалярное произведение как
$n$ умножений и $n-1$ сложений, получится другой счёт; в этом задании
последовательно используем именно соглашение MAC.}

\page{FLOPs: BatchNorm, pooling, head и итог}
В режиме eval BatchNorm использует сохранённые $\mu,\sigma^2$:
\[
 y=\gamma\frac{x-\mu}{\sqrt{\sigma^2+\varepsilon}}+\beta
   =a x+b,\qquad
 a=\frac{\gamma}{\sqrt{\sigma^2+\varepsilon}},\quad b=\beta-a\mu.
\]
В нашей модели считаем 1 умножение и 1 сложение на элемент, то есть
$F_{BN}=2N_{out}$. Подготовку $a,b$ и особенности реализации kernel не считаем.
Это соглашение \texttt{flops()}, а не утверждение, что PyTorch заранее
слил Conv и BN в один оператор.

\begin{center}
\begin{tabular}{lrrrrrrr}
\toprule
 & BN1 & BN2 & BN3 & BN4 & BN5 & BN6 & Сумма \\
\midrule
$N_{out}/X$ & 8 & 4 & 2 & 4 & 1 & 2 & 21 \\
$F_{BN}/X$ & 16 & 8 & 4 & 8 & 2 & 4 & 42 \\
\bottomrule
\end{tabular}
\end{center}

\textbf{ReLU и MaxPool.} Они сравнивают значения. Сравнения не включены
в наш счёт FLOPs, поэтому их вклад здесь нулевой. Это не означает
нулевое время или отсутствие обращений к памяти.

\textbf{GlobalAvgPool.} В одном канале $n=(S/16)^2$ чисел.
Чтобы получить среднее, нужны $n-1$ сложений и 1 деление: всего $n$ операций.
Каналов 512, изображений $B$:
\[
 F_{GAP}=512B\left(\frac{S}{16}\right)^2=2X.
\]
Flatten не выполняет арифметику над значениями.

\textbf{Linear1.} Для каждого из 256 выходов --- 512 MAC и одно сложение bias:
\[
 F_{L1}=B(2\cdot512\cdot256+256)=262400B.
\]
\textbf{Linear2.} Для каждого из 100 выходов --- 256 MAC и bias:
\[
 F_{L2}=B(2\cdot256\cdot100+100)=51300B.
\]
ReLU между Linear также не добавляет FLOPs по принятому соглашению.

\textbf{Общая формула одного forward всего batch:}
\begin{align*}
 F(S,B)&=F_{Conv,total}+F_{BN,total}+F_{GAP}+F_{L1}+F_{L2}\\
 &=17712BS^2+42BS^2+2BS^2+262400B+51300B,\\
 \boxed{F(S,B)}&\boxed{=17756BS^2+313700B\quad\text{FLOPs}.}
\end{align*}
На одно изображение делим на $B$; для GFLOPs делим на $10^9$.

\note{Profiler оценивает прежде всего Conv и Linear. При подсчёте только
их MAC базовая сумма равна $17712BS^2+313344B$. Разница с нашей формулой
составляет $44BS^2+356B$ (BN, GAP и bias). Поэтому небольшое расхождение
с \texttt{flops\_profiler} ожидаемо и не исправляется подбором коэффициента.}

\page{Память модели: каждый параметр и buffer}
Нужно различать \textbf{хранимые данные модели}, \textbf{пиковую память forward}
и \textbf{суммарный трафик}. Сначала найдём постоянную память модели $W$.
Она не умножается на $B$: веса общие для всего batch.

Для Conv без bias число весов $K_i=C_{in,i}C_{out,i}k_i^2$.
У BN на $C$ каналов есть четыре FP32-вектора: обучаемые $\gamma,\beta$
и buffers running mean, running variance. Ещё есть отдельный int64-счётчик
\texttt{num\_batches\_tracked} размером 8 bytes. Он хранится и в eval,
хотя статистика при forward не обновляется.

\begin{center}\small
\begin{tabular}{lrrr}
\toprule
Блок & Веса Conv $K_i$ & BN: FP32-чисел $4C$ & Всего bytes: $4K_i+16C+8$ \\
\midrule
Conv1+BN1 & $3\cdot32\cdot49=4704$ & 128 & 19\,336 \\
Conv2+BN2 & $32\cdot64\cdot25=51200$ & 256 & 205\,832 \\
Conv3+BN3 & $64\cdot128\cdot9=73728$ & 512 & 296\,968 \\
Conv4+BN4 & $128\cdot256=32768$ & 1024 & 135\,176 \\
Conv5+BN5 & $256\cdot256\cdot9=589824$ & 1024 & 2\,363\,400 \\
Conv6+BN6 & $256\cdot512=131072$ & 2048 & 532\,488 \\
\midrule
Сумма & 883\,296 & 4992 & 3\,553\,200 \\
\bottomrule
\end{tabular}\end{center}

Linear1 хранит $512\cdot256+256=131328$ FP32-чисел:
$525312$ bytes. Linear2 хранит $256\cdot100+100=25700$ чисел:
$102800$ bytes. У pooling, ReLU и Flatten параметров нет.
\[
 \boxed{W=4(883296+4992+131328+25700)+6\cdot8=4181312\bytes.}
\]

\textbf{Почему v2 использует $W_{alloc}$, а не просто $W$?}
В модели native CUDA allocator запрос отдельного тензора округляем вверх
до 512 bytes [2]. Для размера $n$ bytes:
\[
 r(n)=512\left\lceil\frac{n}{512}\right\rceil,\qquad
 W_{alloc}=\sum_{\text{отдельные тензоры }j}r(n_j).
\]
Нельзя сначала сложить размеры, а затем округлить один общий результат.
Например, в BN1 каждый из четырёх векторов имеет $32\cdot4=128$ bytes,
каждый округляется до 512; счётчик $8$ тоже до 512.
Вес Conv1: $18816\to18944$. Блок1: $18944+4\cdot512+512=21504$.

\begin{center}\small
\begin{tabular}{lrrrrrrrr}
\toprule
Блок & 1 & 2 & 3 & 4 & 5 & 6 & Linear1 & Linear2 \\
\midrule
$W_{alloc}$ & 21504 & 207360 & 297472 & 135680 & 2363904 & 532992 & 525312 & 102912 \\
\bottomrule
\end{tabular}\end{center}
\[
 \boxed{W_{alloc}=4187136\bytes,\qquad W_{alloc}-W=5824\bytes.}
\]
Это расчёт по формам тензоров при стандартном округлении. Повторное
использование более крупных blocks и нестандартные настройки allocator
могут увеличить фактически занятую память сверх этой оценки.

\page{Пик активаций: почему получается $76BS^2$}
Метрика задания --- \texttt{torch.cuda.max\_memory\_allocated()}.
Она отвечает на вопрос: сколько памяти через allocator было занято
\emph{одновременно} в самый дорогой момент? Складывать выходы всех слоёв нельзя:
в inference старые промежуточные тензоры обычно освобождаются.

Вход $x$ остаётся живым во внешнем коде измерения на протяжении всего forward:
$3X$ чисел. Для обычного оператора в момент создания выхода одновременно
живы его вход и выход. Для BN это два разных массива; для inplace ReLU
выход использует тот же storage. В таблице пока учитываем только основные
FP32-тензоры, без библиотечных временных buffers и постоянной памяти модели.

\begin{center}
\begin{tabular}{lrl}
\toprule
Момент & Живых FP32-чисел & Почему \\
\midrule
Conv1 & $11X$ & $3X+8X$: исходный вход считаем один раз \\
BN1 & $19X$ & $3X+8X+8X$ \\
ReLU1 & $11X$ & $3X+8X$, inplace \\
MaxPool & $13X$ & $3X+8X+2X$ \\
Conv2 & $9X$ & $3X+2X+4X$ \\
BN2 & $11X$ & $3X+4X+4X$ \\
Conv3 & $9X$ & $3X+4X+2X$ \\
BN3 & $7X$ & $3X+2X+2X$ \\
Conv4 & $9X$ & $3X+2X+4X$ \\
BN4 & $11X$ & $3X+4X+4X$ \\
Conv5 & $8X$ & $3X+4X+X$ \\
BN5 & $5X$ & $3X+X+X$ \\
Conv6 & $6X$ & $3X+X+2X$ \\
BN6 & $7X$ & $3X+2X+2X$ \\
\bottomrule
\end{tabular}
\end{center}
ReLU2--6 не создают новую копию активации, поэтому пика соответствующего
BN не превышают.

\textbf{Проверка head.} Даже если одновременно сохранить исходный вход
($3X$), выход features ($2X$), GAP ($512B$), Linear1 ($256B$)
и Linear2 ($100B$), получится лишь $5X+868B$.
Это заведомо щедрая верхняя оценка основных тензоров head; view Flatten
и inplace ReLU не добавляем. При $S\ge32$:
\[
 5X+868B=\left(5+\frac{868}{S^2}\right)X\le5.848X<19X.
\]
Поэтому максимум \emph{этой модели активаций} достигается на BN1:
\[
 M_{activ}=4\max_t N_{live}(t)=4\cdot19X=76BS^2\bytes.
\]
В развёрнутом виде: $12BS^2$ bytes исходного входа плюс два массива по
$32BS^2$ bytes на BN1. При $S$ кратном 16 все три размера кратны 512.

\note{Таблица не описывает скрытые workspace, служебные тензоры kernels
или возможные индексы pooling. Она доказывает максимум для выбранной
модели основных активаций, а не точный peak любого CUDA backend.}

\page{Память v2: постоянный workspace и границы точности}
Одних весов и основных активаций недостаточно. После прогрева Linear
библиотека cuBLAS может сохранять рабочий buffer, выделенный через CUDA
allocator PyTorch. В реализации PyTorch 2.11 стандартный размер для T4
равен [1]:
\[
 W_{BLAS}=(2\cdot4096+8\cdot16)\cdot1024
          =8519680\bytes=8.125\,\mathrm{MiB}.
\]
Это сумма двух blocks по 4096 KiB и восьми по 16 KiB; она взята из
\texttt{parseChosenWorkspaceSize()}, а не подобрана по измеренным ошибкам.

\textbf{Почему он прибавляется к раннему пику BN1, хотя Linear стоит в конце?}
Сначала выполняются три прогревочных forward. Buffer, созданный при Linear
на прогреве, остаётся живым. Затем сбрасывается только счётчик peak;
\texttt{reset\_peak\_memory\_stats()} не освобождает живые allocations.
На следующем forward этот buffer уже есть во время BN1.

Таким образом, текущая функция \texttt{memory()} задаёт:
\begin{align*}
 M(S,B)&=\underbrace{W_{alloc}}_{\text{модель}}+\underbrace{W_{BLAS}}_{\text{workspace}}+\underbrace{M_{activ}}_{\text{живые активации}},\\
       &=4187136+8519680+76BS^2,\\
 \boxed{M(S,B)}&\boxed{=12706816+76BS^2\quad\text{bytes}.}
\end{align*}
Первые два слагаемых постоянны: их сумма $12706816$ bytes. Третье
слагаемое зависит от batch и размера изображения: $76BS^2$ bytes.
Ни один коэффициент в этой формуле не калибруется. В сравнении со старой
формулой $M_{v1}=4181312+76BS^2$ добавка составляет
\[
 M_{v2}-M_{v1}=8525504\bytes.
\]

\textbf{Условия применимости.} Прогретая T4, FP32, стандартный native
allocator, один используемый поток CUDA/рабочий handle, стандартный cuBLAS
workspace, отсутствие дополнительных живых CUDA-тензоров пользователя.
Не задаются overrides размеров workspace или округления allocations.

\textbf{Что остаётся вне формулы.} Дополнительная память kernels cuDNN,
возможный отдельный cuBLASLt workspace, временные служебные тензоры и
особенности повторного использования blocks. Эти расходы зависят от
backend, формы и выбранного алгоритма. Поэтому математически точно
вычислена \emph{наша оценка}, но она не равна гарантированно измеренному пику.

\textbf{Что не нужно просто прибавлять.} CUDA context и свободные cached
blocks относятся к другим расходам устройства; свободный reserved storage
не является живой allocated memory. Сумма тензоров всех слоёв тоже неверна.

\textbf{Оценка OOM.} Если условно доступно $C$ bytes, формула даёт границу
\[
 B_{max}^{model}(S)=\left\lfloor\frac{C-12706816}{76S^2}\right\rfloor
 \quad(C>12706816).
\]
Она не гарантирует отсутствие OOM: физическая память расходуется и на то,
чего нет в модели. Для записанной ёмкости устройства $C=15637086208$ и
$S=512$ эта оценка даёт $B_{max}^{model}=784$; экстраполяция за измеренную
сетку здесь не проверена. В 132 фактических точках OOM не было.

\page{Bytes moved: расчёт трафика для каждого слоя}
$Q(S,B)$ --- суммарный объём чтения и записи за forward. Это не пиковая
занятая память: один storage можно читать и перезаписывать много раз.
В модели каждую активацию читаем/пишем один раз на оператор; веса Conv
читаем один раз на forward. Повторное чтение из DRAM, cache reuse,
fusion, временные buffers и служебные данные не моделируем.

Пусть $A_{in},A_{out}$ --- числа FP32-элементов, $K=C_{in}C_{out}k^2$.
До перевода в bytes получаем:
\begin{align*}
 q_{Conv}&=A_{in}+K+A_{out},\\
 q_{BN}&=A_{out}+4C_{out}+A_{out},\\
 q_{ReLU}&=A_{out}+A_{out},\\
 \boxed{q_{block}}&\boxed{=A_{in}+5A_{out}+K+4C_{out}.}
\end{align*}
Inplace экономит storage, но ReLU всё равно читает и записывает значения.
Четыре BN-вектора читаются один раз на блок, не на каждый пиксель.

\begin{center}\small
\begin{tabular}{lrrrrl}
\toprule
Блок & $A_{in}$ & $A_{out}$ & $K$ & $4C_{out}$ & $q_{block}$ \\
\midrule
1 & $3X$ & $8X$ & 4704 & 128 & $43X+4832$ \\
2 & $2X$ & $4X$ & 51200 & 256 & $22X+51456$ \\
3 & $4X$ & $2X$ & 73728 & 512 & $14X+74240$ \\
4 & $2X$ & $4X$ & 32768 & 1024 & $22X+33792$ \\
5 & $4X$ & $X$ & 589824 & 1024 & $9X+590848$ \\
6 & $X$ & $2X$ & 131072 & 2048 & $11X+133120$ \\
\midrule
Сумма & & & & & $121X+888288$ \\
\bottomrule
\end{tabular}\end{center}

Оставшиеся операторы (всё ещё в FP32-элементах):
\begin{align*}
 q_{MaxPool}&=8X+2X=10X,\\
 q_{GAP}&=2X+512B,\qquad q_{Flatten}=0,\\
 q_{L1}&=512B+512\cdot256+256B=768B+131072,\\
 q_{ReLU,head}&=2\cdot256B=512B,\\
 q_{L2}&=256B+256\cdot100+100B=356B+25600.
\end{align*}
Суммируем и умножаем на 4 bytes/FP32:
\begin{align*}
 Q&=4\bigl[(121+10+2)X+(512+768+512+356)B\\
 &\hspace{38mm}+888288+131072+25600\bigr],\\
 \boxed{Q(S,B)}&\boxed{=532BS^2+8592B+4179840\quad\text{bytes}.}
\end{align*}
\textbf{Точное соответствие коду.} Текущая \texttt{bytes\_moved()} опускает
чтение bias Linear (всего $4(256+100)=1424$ bytes за forward при чтении
по одному разу). В storage и FLOPs bias учтён. Здесь сохраняем именно
реализованную оценку трафика, а не выдаём её за полный аппаратный счётчик.

\page{Latency: смысл каждого члена и два режима v2}
Измеренное время --- медиана семи forward после трёх прогревов; перед
и после каждого измерения выполняется CUDA synchronization.
Время одного forward складывается из накладных расходов и основной работы.
Обозначим: $t_0$ --- совокупные накладные расходы одного forward, seconds;
$P$ --- эффективная скорость арифметики, FLOPs/s;
$R$ --- эффективная скорость передачи данных, bytes/s.

Если выполнить $F$ операций со скоростью $P$, нужно $F/P$ seconds.
Передать $Q$ bytes со скоростью $R$ --- $Q/R$ seconds. В roofline-приближении
арифметика и обмен памятью частично перекрываются, поэтому используем максимум:
\[
 T=t_0+\max\left(\frac{F}{P},\frac{Q}{R}\right).
\]
Это агрегатная модель всей последовательной сети. Она не утверждает полного
перекрытия всех операторов; $t_0$ также не является временем одного kernel.

В v2 допускаем изменение эффективной скорости при большом batch:
\[
 P(B)=\begin{cases}
 P_1,&B\le B_*,\\
 P_2,&B>B_*,
 \end{cases}\qquad P_2\le P_1.
\]
Порог и необходимость второго режима выбираются по train. Это эмпирическое
описание смены эффективности; по одному fit нельзя доказать, какой именно
CUDA kernel или аппаратный эффект её вызвал.

\textbf{Полная формула с раскрытыми $F$ и $Q$:}
\[
\boxed{T(S,B)=t_0+\max\left(
 \frac{17756BS^2+313700B}{P(B)},
 \frac{532BS^2+8592B+4179840}{R}\right).}
\]
В последнем прогоне (коэффициенты ниже округлены для переписывания):
\begin{align*}
 t_0&\approx0.000732280281\ \text{s}=0.732280281\ \text{ms},\\
 P_1&\approx3.7150824604\cdot10^{12}\ \text{FLOPs/s},\\
 P_2&\approx1.9394496515\cdot10^{12}\ \text{FLOPs/s},\\
 R&\approx1.1999933598\cdot10^{11}\ \text{bytes/s},\qquad B_*=64.
\end{align*}
Полная точность сохранена в \texttt{results\_v2/theta.json}.

\textbf{Как читать режимы.} Сравниваем три времени: $t_0$, $F/P(B)$, $Q/R$.
Если больше всего $t_0$, модель launch-bound; если $Q/R$ --- memory-bound;
если $F/P(B)$ --- compute-bound. При близких членах режим смешанный.
Арифметическая интенсивность $I=F/Q$ (FLOPs/byte): между compute и traffic
граница $I=P(B)/R$, если накладные расходы не доминируют.

\note{В последнем прогоне модель относит 98 точек к compute-bound и
34 к launch-bound; memory-bound точек нет. Поэтому fitted $R$ слабо
определяется данными. Эти $P,R$ нельзя интерпретировать как паспортные
характеристики T4. Между измеренными $B=64$ и $B=95$ точное положение
физической границы смены режима неизвестно.}

\page{Как получены коэффициенты latency}
Всего 132 пары $(S,B)$: 11 размеров изображения и 12 batch sizes.
Seed 2026 заранее выделяет 106 train и 26 validation. Из fit исключаются
OOM, нечисловые и неположительные измерения. Параметры и порог выбираются
по train; validation используется для итогового отчёта.

Для точки $i$ обозначим измеренное время $t_i$, а
$f_i=F(S_i,B_i)/10^9$, $q_i=Q(S_i,B_i)/10^9$.
Вместо прямой оптимизации скоростей используем обратные скорости:
\[
 a=\frac{10^9}{P_1},\qquad c=\frac{10^9}{R},\qquad
 d=\frac{10^9}{P_2}-\frac{10^9}{P_1}\ge0.
\]
Тогда для фиксированного порога $b_*$:
\[
 \widehat t_i=t_0+\max\bigl(f_i(a+d\,\mathbf1_{B_i>b_*}),\ c q_i\bigr).
\]
$\mathbf1$ равно 1 для большого batch и 0 иначе. Параметры
$t_0,a,c,d\ge0$; неотрицательность $d$ задаёт $P_2\le P_1$.
При одном режиме дополнительный параметр $d$ отсутствует.

\textbf{Целевая функция fit.} Минимизируем сумму квадратов масштабированных
остатков:
\[
 \min_{t_0,a,c,d\ge0}\sum_{i\in\mathcal T}
 \left(\frac{\widehat t_i-t_i}{s_i}\right)^2,\qquad
 s_i=\max\left(t_i,\ 0.05\,\operatorname{median}_{j\in\mathcal T} t_j\right).
\]
$\mathcal T$ --- только текущая обучающая часть. Нормировка не даёт большим
временам полностью подавить малые. Нижняя граница масштаба защищает от
чрезмерного веса очень маленьких значений. Это weighted least squares,
а не минимизация median APE.

\textbf{Выбор одного/двух режимов по cross-validation:}
\begin{enumerate}
\item Берём baseline с одним $P$ и кандидаты порога из train batch sizes.
По обе стороны каждого порога должно оказаться не меньше 20\% train-точек.
\item Делим только train на пять folds. Для каждого порога пять раз
подгоняем параметры по четырём folds и предсказываем оставшийся.
\item По объединённым отложенным предсказаниям считаем mean APE:
$\frac{100}{N}\sum_i|\widehat t_i-t_i|/t_i$.
\item Выбираем лучший mean APE. Два режима принимаем, только если
$\mathrm{CV}_{two}<0.95\,\mathrm{CV}_{one}$; при равенстве код оставляет baseline.
\item Заново подгоняем выбранную модель на всех 106 train-точках.
\end{enumerate}
В каждом fit пробуем девять начальных приближений и выбираем успешное
с минимальной суммой квадратов. Это уменьшает зависимость от старта,
но не является доказательством глобального оптимума.

В сохранённом прогоне baseline CV mean APE $25.3538\%$, при $B_*=64$
--- $12.9261\%$. Улучшение около $49.0\%$ относительно baseline,
поэтому второй режим принят. Эти CV-ошибки не являются ошибками на
26 итоговых validation-точках.

\page{Energy: мощность, работа и калибровка}
Энергия --- интеграл мощности по времени. Для почти равномерных отсчётов
NVML в серии из $N$ forward используем:
\[
 E^{meas}\approx\frac{\overline{\mathcal P}_{GPU}\,\Delta t}{N}.
\]
Здесь $\overline{\mathcal P}_{GPU}$ --- средняя измеренная мощность всей GPU
в watts, $\Delta t$ --- время серии в seconds; результат в joules на forward.
Это энергия всего batch, не одного изображения. Измеренная мощность уже
содержит базовое потребление GPU, вычитать idle из итогового измерения нельзя.

Модель v2:
\[
\boxed{E(S,B)=(P_{idle}+P_{active})T(S,B)
 +e_0+e_F\frac{F(S,B)}{10^9}+e_Q\frac{Q(S,B)}{10^9}.}
\]
$P_{idle}$ --- измеренная idle power (W); $P_{active}$ --- дополнительный
подгоняемый член мощности (W); $e_0$ --- фиксированные joules/forward;
$e_F$ --- J/GFLOP, $e_Q$ --- J/GB. Здесь $P_{idle},P_{active}$ обозначают
мощность, а $P(B)$ из latency --- производительность: это разные величины.

Член $P_{active}T$ позволяет описать зависящие от длительности расходы
активной GPU сверх idle. $e_F$ и $e_Q$ описывают дополнительные расходы,
связанные с модельной работой. Это совместная регрессионная модель;
по её коэффициентам нельзя строго разложить физическую энергию по узлам GPU.

\textbf{Откуда берутся коэффициенты?} Сначала фиксируем fitted latency.
Берём $P_{idle}$ как медиану idle measurements по допустимым train-строкам.
Для каждой train-точки:
\[
 T_i=T(S_i,B_i;\widehat\theta),\quad
 y_i=E_i^{meas}-P_{idle}T_i,\quad
 z_i=(1,\ F_i/10^9,\ Q_i/10^9,\ T_i).
\]
Ищем $\beta=(e_0,e_F,e_Q,P_{active})\ge0$ методом NNLS:
\[
 \min_{\beta\ge0}\sum_{i\in\mathcal T}
 \left(\frac{z_i\beta-y_i}{u_i}\right)^2,\qquad
 u_i=\max(E_i^{meas},\operatorname{quantile}_{0.05,\mathcal T}E^{meas}).
\]
Используем \emph{предсказанное} $T_i$, чтобы итоговая функция зависела
только от $(S,B,\theta)$. Масштабы и квантиль вычисляются только по train.

В последнем прогоне:
\[
 P_{idle}=30.97,\quad P_{active}\approx27.7825483463,\quad
 e_0=0,\quad e_F\approx0.00495013785384,\quad e_Q=0.
\]
Следовательно, для этих коэффициентов:
\[
 \boxed{E(S,B)\approx58.7525483463\,T(S,B)
 +0.00495013785384\,\frac{17756BS^2+313700B}{10^9}\quad\text{J}.}
\]
Нулевой $e_Q$ не означает бесплатный обмен памятью: в данных его стоимость
не отделилась от остальных членов. NVML имеет ограниченную частоту обновления;
серия для energy и одиночные synchronized forwards для latency могут давать
разные накладные расходы.

\page{Численный пример и проверка точности}
Возьмём $S=224$, $B=8$; тогда $X=8\cdot224^2=401408$.
Всё ниже --- \emph{предсказания}, полученные подстановкой в формулы.
\begin{align*}
 F&=17756\cdot401408+313700\cdot8\\
  &=7129910048\ \text{FLOPs}=7.129910048\ \text{GFLOPs},\\
 Q&=532\cdot401408+8592\cdot8+4179840\\
  &=217797632\bytes,\\
 M&=12706816+76\cdot401408=43213824\bytes\\
  &=41.2119140625\ \mathrm{MiB}.
\end{align*}
Поскольку $8\le64$, используем $P_1$. Из коэффициентов полной точности:
\begin{align*}
 F/P_1&\approx0.001919179486\ \text{s},\\
 Q/R&\approx0.001814990310\ \text{s},\\
 T&\approx0.000732280281+0.001919179486\\
  &\approx0.002651459767\ \text{s}=2.651459767\ \text{ms},\\
 E&\approx58.7525483463\cdot0.002651459767
     +0.00495013785384\cdot7.129910048\\
  &\approx0.191074056\ \text{J}.
\end{align*}
Вторая контрольная точка проверяет ветку большого batch:
\begin{center}\small
\begin{tabular}{lrr}
\toprule
 & $S=32,\ B=1$ & $S=224,\ B=128$ \\
\midrule
$F$, FLOPs & 18\,495\,844 & 114\,078\,560\,768 \\
$Q$, bytes & 4\,733\,200 & 3\,422\,064\,512 \\
$M$, bytes & 12\,784\,640 & 500\,818\,944 \\
$T$, ms & 0.771724 & 59.552348 \\
$E$, J & 0.0454323 & 4.0635568 \\
Ветка & $P_1$ & $P_2$ \\
\bottomrule
\end{tabular}\end{center}

\textbf{Как оценена точность.} Для положительного измерения $y_i$:
$\mathrm{APE}_i=100|\widehat y_i-y_i|/y_i$.
Median APE --- медиана этих индивидуальных ошибок, mean APE --- их среднее.
Это не ошибка медианы измерений и не одна ошибка общего среднего.

\begin{center}\small
\begin{tabular}{lrrrr}
\toprule
Validation, 26 точек & v1 median & v2 median & v1 mean & v2 mean \\
\midrule
FLOPs & $0.25\%$ & $0.25\%$ & $0.25\%$ & $0.25\%$ \\
Memory & $28.90\%$ & $22.03\%$ & $35.75\%$ & $21.90\%$ \\
Latency & $27.83\%$ & $14.63\%$ & $28.26\%$ & $16.21\%$ \\
Energy & $21.00\%$ & $6.89\%$ & $24.09\%$ & $11.13\%$ \\
\bottomrule
\end{tabular}\end{center}
Grid и split совпадают, но прогоны разные. Все 132 измерения памяти
совпали, поэтому её улучшение связано с формулой. Сравнение latency/energy
также включает вариативность GPU. Holdout уже использовался в прежних
отчётах; для независимой проверки новой гипотезы нужны свежие конфигурации.

\page{Итог для рукописной работы и источники}
Четыре функции для этой сети, FP32, eval, inference, $S\in16\mathbb N$,
$S\ge32$ и положительного целого $B$:
\begin{align*}
 \boxed{F(S,B)}&\boxed{=17756BS^2+313700B,}\\[4pt]
 \boxed{M(S,B)}&\boxed{=12706816+76BS^2,}\\[4pt]
 \boxed{T(S,B;\theta)}&\boxed{=t_0+\max(F(S,B)/P(B),\ Q(S,B)/R),}\\[4pt]
 \boxed{E(S,B;\theta_E)}&\boxed{=(P_{idle}+P_{active})T+e_0+e_FF/10^9+e_QQ/10^9.}
\end{align*}
Вспомогательные определения:
\[
 Q(S,B)=532BS^2+8592B+4179840,\qquad
 P(B)=\begin{cases}P_1,&B\le B_*\\P_2,&B>B_*\end{cases}.
\]
$F$ измеряется в FLOPs, $M,Q$ --- в bytes, $T$ --- в seconds, $E$ --- в joules.
Значения относятся к одному forward всего batch. Только latency и energy
содержат fitted parameters. Все суммы FLOPs, storage и модельного трафика
получены из размеров слоёв; workspace памяти --- из настройки реализации.

\textbf{Что переписать, чтобы вывод оставался полным.}
Сначала соглашения и таблицу форм. Затем общую формулу Conv и подстановки
для шести слоёв; BN, GAP, head и сумму FLOPs. Для памяти --- таблицу весов,
округление allocations, максимум одновременно живых активаций и workspace.
После этого --- таблицу traffic, latency с двумя режимами, energy и обе
целевые функции fit. В конце --- один численный пример и ограничения.
Страницы этой заготовки соответствуют этим шагам; при необходимости
таблицы можно разнести на несколько рукописных листов.

\textbf{Что обязательно сохранить в объяснениях.}
Разделять $M$ и $Q$; не умножать веса на batch; не считать inplace ReLU
бесплатным по трафику; не называть оценку памяти точным измерением.
Объяснить, почему warmup workspace присутствует уже при BN1. Указать,
что $P(B)$ и $R$ --- эффективные параметры, а выбор порога выполнен внутри
train. Наблюдаемое улучшение ошибки не снимает ограничений модели.

\begingroup\raggedright
\textbf{Основа и воспроизводимость.}
Архитектура: \path{hw1/models.py}. Формулы: \path{hw1/equations.py}.
Калибровка: \path{hw1/calibrate.py}. Протокол: \path{hw1/measure.py}.
Коэффициенты и ошибки: \path{hw1/results_v2/theta.json},
\texttt{errors.json}; сравнение: \path{hw1/versions_comparison.json}.
Текущий автономный ноутбук: \path{hw1/hw1_version_2.ipynb}.

\par\endgroup
\textbf{Первичные источники для констант реализации.}

[1] PyTorch v2.11.0, \texttt{CublasHandlePool.cpp}:
\texttt{parseChosenWorkspaceSize()}, \texttt{getNewWorkspace()},
\texttt{setWorkspaceForHandle()}.
\href{https://github.com/pytorch/pytorch/blob/v2.11.0/aten/src/ATen/cuda/CublasHandlePool.cpp}{Исходный код cuBLAS workspace}.

[2] PyTorch v2.11.0, \texttt{CUDACachingAllocator.cpp},
\texttt{round\_size()}: стандартное округление запроса.
\href{https://github.com/pytorch/pytorch/blob/v2.11.0/c10/cuda/CUDACachingAllocator.cpp}{Исходный код CUDA allocator}.

Это печатная заготовка для понимания и переписывания. Итоговый рукописный
PDF пользователь составляет самостоятельно; первоначальный
\texttt{hw1\_handwritten.pdf} этим файлом не заменён.
\end{document}
